% 03-core.tex: the credential core.
\section{The core: one credential, one signature, one person}
\label{sec:core}

\subsection{What the credential is}

At the centre of Polaris is one row of the table \entity{IdentityToken} and at least one row of
\entity{TokenSignature}. The token row records a unique \code{token\_value}, a physical serial, the
person it belongs to, the agency that issued it, the algorithm it was issued under, its predecessor if
it replaced an earlier token, and its status. The signature row records the issuing agency's ML-DSA
signature over the SHA3-256 digest of the token value, together with the public key that produced
it, so that verification is self-contained and survives the retirement of the signing key. A token
without a signature row cannot exist: a database trigger refuses any write that would leave one
without a non-deprecated signature. The credential is therefore not a document but a signed claim:
\emph{this authority issued this token to this person.} \sImpl

What the holder holds is the authenticity pack: the token value, the signature and the issuer's
public key, exported by the application and kept as a file by the reference wallet
(\code{scripts/polaris-wallet.py}). Anyone can check the pack with a standard ML-DSA library and no
Polaris code (Section~\ref{sec:verifiers}). The physical token that Version 1 modelled, a hardware
card with on-device biometric binding, remains modelled and not manufactured: the schema carries the
serial, the binding type and the device bindings, and no card exists. \sFut{} for the hardware;
\sImpl{} for the record.

\figfile{nucleus}{The core entity geometry, unchanged from Version 1. Seven entities around
\entity{IdentityToken}: the person who holds it, the agency that issues it, the algorithm that signs
it, the contexts it may be verified in, and the two append-only event logs it generates. The
many-to-many relationships resolve to the junction tables \entity{AgencyAlgorithmAuth} and
\entity{TokenPermission}. This geometry is the constant of the paper: Section~\ref{sec:walls} draws
the other twenty-nine tables as rings around it.}{fig:nucleus}

\subsection{What the credential is not}

The constitution's tenth constraint, C10, states that identity attestation never carries spending
authority: there is no monetary claim in the schema, and a check fails the build if one appears. The
reason is stated as a failure mode, not a preference. A credential that is also money inherits
programmability: every constraint anyone wants to place on spending accretes onto the identity token,
until an administrative paperwork error becomes an existential bank-balance error, and until the
system can be told that a person may not buy fuel. The same refusal covers reputation, behaviour and
profiles. Polaris answers one question, \emph{is this token valid for this context right now}, and
it declines to answer \emph{should this person be allowed to do X}. That second question belongs to
the relying party, outside the boundary. \sImpl

\subsection{One active token per person}

A person may hold several provisioned tokens but only one may be \code{ACTIVE}. Two active tokens for
one person would let one token authorise a transaction and the other repudiate it. The rule is C3,
and it is enforced by a partial unique index, \code{uq\_one\_active\_per\_person}, on the person's
identifier where the status is active. Two operators activating two reserve tokens for the same
holder at the same moment find that exactly one of them succeeds, and the test that proves it runs
with real threads against a live database (C9). \sImpl

\subsection{The lifecycle}

\figfile{lifecycle}{The token lifecycle machine, unchanged from Version 1. \code{ACTIVE} is the only
operational state. A token enters at \code{RESERVE}, is activated, and leaves through one of three
terminal states or into \code{DORMANT} when a successor takes over. At \code{v9.345} the legal
transitions are enforced by the trigger \code{trg\_token\_state\_machine}, and every transition
writes a lifecycle event automatically.}{fig:lifecycle}

Version 1 recommended enforcing the transitions in a trigger and left it as future work. At HEAD the
trigger exists, and a second trigger writes a \entity{TokenLifecycleEvent} row for every status
change, including a change made by hand from a database shell. Succession is by reference, never by
overwrite: a replacement token points at its predecessor and the old token stays in the database in
its terminal state. Lost tokens stay lost, and their data is not erased. \sImpl

\subsection{Entering, leaving and returning}

Three mechanisms surround the core so that the schema cannot be turned against the holder at the
edges of the lifecycle.

\begin{description}
  \item[Entering: tiered enrolment.] Every person has an enrolment status drawn from five values,
    and \code{EXEMPT} is a first-class one: a recognised way to take part in civic life without a
    token, whether for biometric incompatibility, religious exemption or an alternative path. The
    asymmetry is deliberate. Recording an exemption is one row; enumerating the not-enrolled is
    possible but unhelped by any function, and it is audited when done. Nothing is derived from token
    state, so an administrative revocation never silently becomes a person leaving civic life.
    \sImpl
  \item[Leaving: issuer discretion.] The worst failure of an identity system is not a forged token
    but an authority revoking a population's tokens at scale. A trigger bounds how fast one agency can
    revoke its own tokens (five percent of its outstanding base in thirty days by default) and above
    that a co-signer from a different agency is required, recorded in the audit trail. The bound
    converts a surprise into a trend that reporting can see. It does not stop slow abuse under the
    ceiling, and says so. \sImpl
  \item[Returning: the recovery ceremony.] A fire takes the wallet and the reserve together. Recovery
    is mandatory and is also an impersonation path, so it is made expensive without being made
    impossible: a two-phase ceremony with a 48-hour cool-down, three independent out-of-band channels,
    an approver who cannot be the requester, and admin-only completion, each encoded as a check
    constraint on \entity{RecoveryRequest}. The holder is without identity for the cool-down; the
    substitute credential a production system would want for that window is not built. \sImpl{}
    for the ceremony, \sFut{} for the interim credential.
\end{description}

\subsection{Compulsion: the duress code}

The vocation above the ten constraints is anti-coercion: no person may be compelled to renounce,
transfer or surrender their identity against their will. The mechanism that gives that sentence
something to stand on is old and comes from banking: a second secret that looks exactly like
success to whoever is watching, while an alert goes out through a channel the watcher cannot see.
A holder may enrol a duress code. When it is presented, the verification proceeds identically on
every operator-visible surface, the same page, the same flash message, the same event row, and a
\entity{DuressEvent} row is appended to a table the operator's screens never join, from which a
pager fires at the highest severity. The comparison is constant-time and the work is paid on both
paths. \sImpl{} as a tested property of the modelled flow; \sExt{} as a side-channel guarantee,
because no external party has measured it.

\begin{layercard}{Layer card: the credential core}
\qa{What it is}{One \entity{IdentityToken} row and its \entity{TokenSignature} rows: a signed claim that an authority issued this token to this person, in a lifecycle with one active state.}
\qa{Why it exists}{Every other structure in Polaris organises itself around this claim. It is kept as small as possible so that the rest can be bounded.}
\qa{What it trusts}{The issuing agency's signing key, held behind a custody interface (Section~\ref{sec:signatures}); the schema's constraints.}
\qa{What it refuses}{Application code as a source of truth: the state machine, uniqueness and audit are decided by the database. Money, reputation and profiles: C10.}
\qa{What it can see}{The holder's legal name, date of birth and country; the issuing agency; the algorithm; the lifecycle.}
\qa{What it cannot see}{A biometric template (never stored), a verification graph (Section~\ref{sec:privacy}), a spending history (never modelled).}
\qa{If it fails}{A forged core is a forged signature, which the two-witness verification and the detached verifier refuse; a duplicated core is refused by the unique index; a rewritten core is refused by the append-only triggers.}
\qa{Checked by}{\code{check\_c1c10\_objects\_resolve}, the C3 property tests, the concurrency suite, the attack suite's forgery and tamper cases.}
\qa{Now or later}{Now. It is the part of Polaris that has existed since Version 1.}
\qa{Inside and outside}{Inside: the person, who exists outside the software. Outside: the wall of constraints (Section~\ref{sec:walls}).}
\end{layercard}

\analogy{The town hall keeps one register. Each entry says who a person is, which official entered
it, and under which seal. The register cannot be edited, only appended to; a second entry for the
same person is refused at the desk; and the seal is one that every reader can check against the
official's published mark without asking the hall.}
